% p.69 — الباب الموفي بالثلاثين في العوامل ومتعلقاتها (تتمة)
\mspage{69}
% Layout: two columns (right column first), each a list «رقم + عامل | أمثلة» with the
% particle in a narrow sub-column at the right edge; prose sub-headings («مِنْهَا حُرُوْفٌ …»)
% run across the full width of the column between groups.  The column break falls after
% (31) هَيَا (foot of the right column).  Plain numbers, written as (n).
% Set as `twocol` with `rightbox`/`leftbox` (not rightcol/leftcol: reledpar pairs those
% line by line and a table counts as one line, so every table would open an equally tall
% gap in the other column).  Each box holds its column's sequence of sarf tables and
% \par-separated heading paragraphs; the tables are scaled to the column width.
% Printed page number ۶۹ at top left.
\begin{twocol}
\begin{rightbox}\vspace{0pt}% align the tops of the two boxes (minipage[t] aligns first baselines)
\begin{sarf}{|r|r|}
(17) حَاشَا & جَاءَنِي الْقَوْمُ حَاشَا زَيْدٍ\\
(18) عَدَا & رَأَيْتُ الْقَوْمَ عَدَا زَيْدٍ\\
(19) خَلَا & مَرَرْتُ بِالْقَوْمِ خَلَا زَيْدٍ\\
\end{sarf}
مِنْهَا حُرُوْفٌ تَنْصِبُ الْاِسْمَ وَتَرْفَعُ الْخَبَرَ تُسَمّٰى حُرُوْفًا مُشَبَّهَةً بِالْفِعْلِ وَهِيَ سِتَّةٌ
\par
\begin{sarf}{|r|r|}
(20) اِنَّ & اِنَّ زَيْدًا قَائِمٌ\\
(21) أَنَّ & عَلِمْتُ أَنَّ زَيْدًا قَائِمٌ\\
(22) كَأَنَّ & كَأَنَّ زَيْدًا أَسَدٌ\\
(23) لٰكِنَّ & غَابَ زَيْدٌ لٰكِنَّ بَكْرًا حَاضِرٌ\\
(24) لَيْتَ & لَيْتَ الشَّبَابَ عَائِدٌ\\
(25) لَعَلَّ & لَعَلَّ السُّلْطَانَ حَاضِرٌ\\
\end{sarf}
مِنْهَا حَرْفَانِ تَرْفَعَانِ الْاِسْمَ وَتَنْصِبَانِ الْخَبَرَ تُسَمّٰى مُشَبَّهَتَيْنِ بِلَيْسَ
\par
\begin{sarf}{|r|r|}
(26) مَا & مَا زَيْدٌ قَائِمًا\\
(27) لَا & لَا رَجُلٌ ظَرِيْفًا\\
\end{sarf}
مِنْهَا حُرُوْفٌ تَنْصِبُ الْاِسْمَ فَقَطْ وَهِيَ سَبْعَةٌ
\par
% TODO: the lithograph numbers both أَيَا and هَيَا as ۳۱; the second should be 32.
\begin{sarf}{|r|r|}
(28) اَلْوَاوُ & اِسْتَوَى الْمَاءُ وَالْخَشَبَةَ\\
(29) اِلَّا & جَاءَنِي الْقَوْمُ اِلَّا زَيْدًا\\
(30) يَا & يَا عَبْدَ اللّٰهِ\\
(31) أَيَا & أَيَا غُلَامَ زَيْدٍ\\
\com{(31)}{كذا في الأصل، ولعل الصواب: (32)} هَيَا & هَيَا شَرِيْفَ الْقَوْمِ\\
\end{sarf}
\end{rightbox}
% ---- left column ----
\begin{leftbox}\vspace{0pt}%
\begin{sarf}{|r|r|}
(33) أَيْ & أَيْ أَفْضَلَ الْقَوْمِ\\
(34) أَ & أَ عَبْدَ اللّٰهِ\\
\end{sarf}
مِنْهَا حُرُوْفٌ تَنْصِبُ الْمُضَارِعَ وَهِيَ أَرْبَعَةٌ
\par
\begin{sarf}{|r|r|}
(35) أَنْ & أُرِيْدُ أَنْ أُسْلِمَ\\
(36) لَنْ & لَنْ يَدْخُلَ الْجَنَّةَ مُتَكَبِّرٌ\\
(37) كَيْ & أَسْلَمْتُ كَيْ أَدْخُلَ الْجَنَّةَ\\
(38) اِذَنْ & أَسْلَمْتُ اِذَنْ تَدْخُلَ الْجَنَّةَ\\
\end{sarf}
مِنْهَا حُرُوْفٌ تَجْزِمُ الْمُضَارِعَ وَهِيَ خَمْسَةٌ
\par
\begin{sarf}{|r|r|}
(39) اِنْ & اِنْ تَضْرِبْ أَضْرِبْ . اِنْ تَضْرِبْ ضَرَبْتُكَ\\
 & اِنْ تَضْرِبْ فَزَيْدٌ ضَارِبٌ . اِنْ ضَرَبْتَ أَضْرِبْ\\
(40) لَمْ & لَمْ يَلِدْ وَلَمْ يُوْلَدْ\\
(41) لَمَّا & لَمَّا يَدْخُلِ الْاِيْمَانُ فِيْ قُلُوْبِكُمْ\\
(42) لَامُ الْأَمْرِ & لِيَضْرِبْ زَيْدٌ\\
(43) لَامُ النَّهْيِ & لَا يَضْرِبْ زَيْدٌ لَا تَضْرِبْ\\
\end{sarf}
مِنْهَا أَسْمَاءٌ تَجْزِمُ الْمُضَارِعَيْنِ عَلٰى مَعْنٰى اِنْ وَهِيَ تِسْعَةٌ
\par
\begin{sarf}{|r|r|}
(44) مَنْ & مَنْ يُكْرِمْنِيْ أُكْرِمْهُ\\
(45) مَا & مَا نَنْسَخْ مِنْ آيَةٍ نَأْتِ بِخَيْرٍ مِنْهَا\\
(46) مَتٰى & مَتٰى تَذْهَبْ أَذْهَبْ\\
(47) مَهْمَا & مَهْمَا تَخْرُجْ أَخْرُجْ\\
(48) أَيٌّ & أَيٌّ يُكْرِمْنِيْ أُكْرِمْهُ\\
\end{sarf}
\end{leftbox}
\end{twocol}
